\section{Elipse de acción, radián y doble retorno de holonomía}
\label{sec:elipse-accion-holonomia-rev8}
\index{acción!elipse}
\index{radián!sector de acción}
\index{Möbius!doble retorno}

El teorema focal anterior fija la forma angular y el lector determinantal
fija la escala de acción. Definamos
\[
 \eta=\operatorname{artanh}\frac{C^*}{A},
 \qquad
 a_S=\sqrt{2\hbar_{\rm HMT}^{\rm ret}e^\eta},
 \qquad
 b_S=\sqrt{2\hbar_{\rm HMT}^{\rm ret}e^{-\eta}}.
\]
La órbita simpléctica
\[
 Q(\theta)=a_S\cos\theta,
 \qquad
 P(\theta)=b_S\sin\theta
\]
conserva simultáneamente la razón de forma del bloque angular y el área
fijada por la recta de acción.

\begin{theorem}[Área barrida por fase]
\label{thm:area-fase-elipse-rev8}
Para todo \(\theta\),
\[
 \operatorname{Area}(\theta)
 =\frac12\int_0^\theta(QP'-PQ')\,dt
 =\hbar_{\rm HMT}^{\rm ret}\theta.
\]
Por tanto,
\[
 \operatorname{Area}(1\ {\rm rad})=\hbar_{\rm HMT}^{\rm ret},
 \qquad
 \operatorname{Area}(2\pi)=h_{\rm HMT}^{\rm ret}.
\]
\end{theorem}

\begin{proof}
Como \(QP'-PQ'=a_Sb_S\) y
\(a_Sb_S=2\hbar_{\rm HMT}^{\rm ret}\), la integral vale
\(\hbar_{\rm HMT}^{\rm ret}\theta\). La segunda identidad usa
\(h_{\rm HMT}^{\rm ret}=2\pi\hbar_{\rm HMT}^{\rm ret}\).
\end{proof}

APP aporta el área orientada; \(A,C^*\) fijan la forma y los focos; la rama
determinantal fija la escala. La vuelta primitiva encierra \(h\) y un radián
de fase barre \(\hbar\). El radián ordinario sigue siendo
\(1\,\mathrm{rad}=180^\circ/\pi\); la construcción añade su genealogía como
sector unitario de acción, no otra constante angular.

Para una sección persistente con holonomía
\(\varepsilon_\gamma\in\{+1,-1\}\), la condición
\[
 \varepsilon_\gamma e^{iJ_\gamma/\hbar}=1
\]
distingue dos cierres. En el cilindro, \(\varepsilon_\gamma=+1\) y el ciclo
primitivo transporta \(J=2\pi\hbar=h\). En la banda de Möbius,
\(\varepsilon_\gamma=-1\): una vuelta visible cambia de hoja y transporta
\(J=\pi\hbar=h/2\); la segunda restaura orientación y completa \(h\). Por
radián de la fase interna se transporta \(\hbar\) en ambos casos; el factor
\(1/2\) pertenece a la proyección sobre la vuelta visible de la base Möbius.

\begin{figure}[htbp]
\centering
\resizebox{.98\textwidth}{!}{\begin{tikzpicture}[
  font=\scriptsize,
  box/.style={draw,rounded corners=2mm,align=center,minimum height=.9cm,
    inner xsep=5pt},
  arr/.style={-{Stealth[length=2.2mm]},thick}]

  \begin{scope}[shift={(0,0)}]
    \draw[->,hmtgray] (-3.0,0)--(3.0,0) node[right]{\(u_+\)};
    \draw[->,hmtgray] (0,-2.0)--(0,2.0) node[above]{\(u_-\)};
    \draw[hmtblue,very thick] (0,0) ellipse (2.65cm and 1.45cm);
    \fill[hmtgold] (-1.85,0) circle (2pt) node[below]{\(F_-^{(\lambda)}\)};
    \fill[hmtgold] (1.85,0) circle (2pt) node[below]{\(F_+^{(\lambda)}\)};
    \node[text=hmtblue,font=\bfseries] at (0,2.55) {elipse de tasas};
    \node[align=center,text=hmtgray] at (0,-2.4)
      {\(q_+=e^{-a_\lambda^2}\), \(q_-=e^{-b_\lambda^2}\)};
  \end{scope}

  \draw[arr,hmtgold] (3.3,0)--node[above,align=center]{misma forma\\escala de acción} (6.0,0);

  \begin{scope}[shift={(9.0,0)}]
    \draw[->,hmtgray] (-2.8,0)--(2.8,0) node[right]{\(Q\)};
    \draw[->,hmtgray] (0,-2.0)--(0,2.0) node[above]{\(P\)};
    \fill[hmtlightgold] (0,0)--(2.5,0)
      arc[start angle=0,end angle=57.2958,x radius=2.5cm,y radius=1.35cm]--cycle;
    \draw[hmtviolet,very thick] (0,0) ellipse (2.5cm and 1.35cm);
    \draw[hmtgold,thick] (0,0)--(2.5,0);
    \draw[hmtgold,thick] (0,0)--({2.5*cos(57.2958)},{1.35*sin(57.2958)});
    \node[text=hmtviolet,font=\bfseries] at (0,2.55) {elipse de acción};
    \node[align=center] at (.7,.48) {\(\theta=1\)\\área \(\hbar\)};
    \node[align=center,text=hmtgray] at (0,-2.4) {vuelta \(2\pi\): área \(h\)};
  \end{scope}

  \node[box,draw=hmtblue,fill=hmtlightblue,minimum width=3.2cm] (cyl) at (3.2,-4.3)
    {cilindro \(\varepsilon=+1\)\\\(2\pi\mapsto h\)};
  \node[box,draw=hmtviolet,fill=hmtlightviolet,minimum width=4.2cm] (mob) at (9.0,-4.3)
    {Möbius \(\varepsilon=-1\)\\\(2\pi\mapsto h/2\), \(4\pi\mapsto h\)};
  \draw[arr,hmtgray] (cyl)--node[
    above,align=center,text width=1.9cm,fill=white,inner sep=1pt
  ]{cambio de\\holonomía} (mob);
  \node[align=center,text=hmtgray,font=\footnotesize] at (6.1,-5.55)
    {La fase interna transporta \(\hbar\) por radián;\\el factor dos pertenece a la vuelta visible.};
\end{tikzpicture}
}
\caption{La elipse de tasas fija la forma; su levantamiento a la recta de acción fija
el área. El cierre cilíndrico y el doble retorno de Möbius distinguen la fase
interna de la vuelta visible.}
\label{fig:elipse-accion-holonomia-rev8}
\end{figure}

\subsection{Geometría de Hilbert--Beltrami--Klein en la elipse de acción}
\label{sec:hilbert-beltrami-klein-elipse-accion}

La misma elipse que porta el área simpléctica admite una geometría proyectiva
intrínseca. Sea
\[
 \mathbb E_{\rm act}
 :=\left\{(Q,P)\in\mathbb R^2:
 \frac{Q^2}{a_S^2}+\frac{P^2}{b_S^2}<1\right\},
 \qquad
 \mathcal L_{\rm act}(Q,P)
 :=\left(\frac{Q}{a_S},\frac{P}{b_S}\right).
\]
La aplicación lineal \(\mathcal L_{\rm act}\) identifica
\(\mathbb E_{\rm act}\) con el disco unidad \(\mathbb D\). Para dos puntos
distintos \(x,y\in\mathbb E_{\rm act}\), sean \(a,b\in\partial
\mathbb E_{\rm act}\) las intersecciones de la recta \(xy\) con la frontera,
ordenadas como \(a,x,y,b\). Definimos
\begin{equation}
 d_{\rm H}^{\rm act}(x,y)
 :=\frac12\log
 \frac{\lVert y-a\rVert\,\lVert x-b\rVert}
      {\lVert x-a\rVert\,\lVert y-b\rVert}.
 \label{eq:distancia-hilbert-elipse-accion}
\end{equation}
La expresión es la razón doble de los cuatro puntos colineales y, por ello,
es independiente de la norma auxiliar usada sobre esa recta.

\begin{proposition}[Realización de Beltrami--Klein de la elipse de acción]
\label{prop:beltrami-klein-elipse-accion}
El par \((\mathbb E_{\rm act},d_{\rm H}^{\rm act})\) es isométrico, mediante
\(\mathcal L_{\rm act}\), al modelo de Beltrami--Klein del plano hiperbólico.
Sus geodésicas no orientadas son exactamente las cuerdas abiertas de la
elipse. La frontera \(\partial\mathbb E_{\rm act}\) conserva, de manera
simultánea y tipada, la órbita simpléctica
\[
 \gamma_{\rm act}(\theta)
 =(a_S\cos\theta,b_S\sin\theta),
 \qquad
 \frac12\int_0^\theta(QP'-PQ')\,dt
 =\hbar_{\rm HMT}^{\rm ret}\theta.
\]
\end{proposition}

\begin{proof}
Las transformaciones proyectivas preservan razones dobles. Como
\(\mathcal L_{\rm act}\) es una transformación lineal invertible que lleva
la elipse al disco unidad, transporta
\eqref{eq:distancia-hilbert-elipse-accion} a la distancia de Hilbert del
disco. Ésta es la métrica de Beltrami--Klein, cuyas geodésicas son sus cuerdas.
La identidad simpléctica de la frontera es el
\cref{thm:area-fase-elipse-rev8} escrita en la carta normalizada.
\end{proof}

La carta completa queda, por tanto, tipada por
\[
 \mathfrak R_{\rm act}:
 (A,C^*,\hbar_{\rm HMT}^{\rm ret})
 \longmapsto
 (\mathbb E_{\rm act},d_{\rm H}^{\rm act},
  dP\wedge dQ,\gamma_{\rm act}).
\]
La métrica proyectiva gobierna el interior; la forma simpléctica y la ley de
área gobiernan la frontera recorrida. Una realización como región física o
como dominio de matrices de covarianza se formula mediante un funtor ulterior
que conserve explícitamente estas dos estructuras; la construcción anterior
fija ya su dominio matemático y los invariantes que ese funtor debe preservar.
